\documentclass[10pt,a4paper]{amsart}
\usepackage[margin=19mm]{geometry}
\usepackage{amsmath,amssymb,mathtools}
\usepackage[scr=rsfs,cal=euler]{mathalfa}
\usepackage{xfrac}
\usepackage[unicode,hidelinks,bookmarks=false]{hyperref}
\newcommand{\ie}{i.e.\ }
\newcommand{\cf}{cf.\ }
\newcommand{\resp}{respectively\ }
\newcommand{\Subsection}{\subsection}
\providecommand{\nobreakdash}{\nobreak\hbox{-}}
\setlength{\emergencystretch}{2em}
\multlinegap0pt
\usepackage{tikz}
\usepackage{bm}
\usepackage{amssymb}
\usepackage{mathtools}
\usepackage[shortlabels]{enumitem}
\newtheorem{proposition}{Proposition}
\newtheorem{lemma}{Lemma}
\newtheorem{theorem}{Theorem}
\usepackage{cleveref}
\crefname{proposition}{Proposition}{Propositions}
\crefname{lemma}{Lemma}{Lemmas}
\crefname{theorem}{Theorem}{Theorems}
\title{Lagrangian correspondences for moduli spaces of Higgs bundles and holomorphic connections}
\author{Panagiotis Dimakis, Duong Dinh, Shengjing Xu}
\date{Source: arXiv:2604.14127v1 (CC BY 4.0)}
\begin{document}
\maketitle

\noindent\emph{Source and licence.} Lagrangian correspondences for moduli spaces of Higgs bundles and holomorphic connections. Version arXiv:2604.14127v1 (CC BY 4.0), distributed by arXiv under the Creative Commons Attribution 4.0 International licence, http://creativecommons.org/licenses/by/4.0/, as recorded in the arXiv licence metadata for this exact version and shown on the abstract page https://arxiv.org/abs/2604.14127v1. Full licence notice: \texttt{licenses/arxiv-2604.14127-CC-BY-4.0.txt}.\par\smallskip

\noindent\textit{Editorial scope.} Counted unit: the article's theorem SLn (source0.tex lines 3134--3152). The article's complete original proof of it is delivered with it (source0.tex lines 3272--3430, one \texttt{proof} environment of 159 lines, which the article places away from the statement). No mathematics is written, repaired or re-derived: the statement and the proof are the article's own lines, in the article's own notation, and no retained line is substituted or re-flowed.  Retained uncounted as the source-local context the proof needs: lines 1991--1996; lines 2060--2071; lines 2280--2295; lines 2497--2519; lines 2623--2638; lines 3155--3176.  Nothing was omitted from any retained span, so the package carries no omission record.  No retained line contains a citation, so the wrapper carries no bibliography and no bibliography entry.  No retained line contains a figure environment, an image-inclusion command or drawing code, so no image is delivered and the package declares no asset dependency.  Editorial formatting only, in addition: an AMS \texttt{amsart} preamble that restates the article's own declarations and macros used by the retained lines, namely the environments \texttt{proposition}, \texttt{lemma}, \texttt{theorem} on separate counters, together with the standard packages those lines need.  The retained environments print in this document as Proposition 1, Lemma 1, Theorem 1 in order of appearance, whereas the article prints the same items with the numbering of its own full document.  The complete native source of the article as distributed in the arXiv e-print for this exact version is bundled unchanged as \texttt{sources/198520/source0.tex}.
\begin{proposition}\label{p:oper bundle out D}
There is a canonical isomorphism 
\(
\mathcal F_i^\vee \;\cong\; J^{i-1}(L^{-1}).
\)
\end{proposition}

\begin{equation}\label{eqn-nabla-filtration}
    \nabla = \mathrm{d} +
    \begin{pmatrix}
        &0  &\ast &\dots &\ast &\ast &\ast  \\
        &1 &0 &\dots &\ast &\ast &\ast \\
        &0 &1 &\dots &\ast &\ast &\ast \\
        &\vdots  &\vdots &\dots &\vdots  &\vdots &\vdots \\
        &0  &0 &\dots &0 &\ast &\ast \\
        &0  &0 &\dots &1 &\ast &\ast \\
        &0  &0 &\dots &0 &s_i(\nabla) &\ast 
    \end{pmatrix}.
\end{equation}

\begin{proposition}\label{p: generic finite dR}
Fix a line bundle \(L\) such that
\(
L^{\otimes n}\simeq K_C^{\otimes \binom{n}{2}}(-D).
\)
Consider the natural forgetful morphism
\[
\pi(L,D)\colon
\widetilde{\mathbb{L}_{dR}}(L,D)
\longrightarrow
\mathbb{L}_{dR}(L,D),
\qquad
([i]\colon L\hookrightarrow E,\nabla)\longmapsto (E,\nabla),
\]
then for generic \(D\in \operatorname{Hilb}^d(C)\), the morphism \(\pi(L,D)\) is generically finite.
\end{proposition}

\begin{lemma}\label{p:trans of diff opers}
View $\mathbf D_m$ (resp.\ $\mathbf D$) as a collection of compatible local scalar differential operators which patch to a global operator on $C\setminus D$ (resp.\ on $C$). Fix the above local charts. Then the following hold.
\begin{itemize}
    \item 
$\mathbf D_m$ and $\mathbf D$ agree on $C\setminus D$.
Over $U_i$, one has
\begin{equation}\label{e:transform rule of mero oper0}
\mathbf D(z)
=
z^{\frac{n+1}{n}}\circ \mathbf D_m(z) \circ z^{-\frac{1}{n}} .
\end{equation}

\item Locally, write
\begin{equation}\label{local Dm}
\mathbf D_m
:= \partial^n + Q_2(z)\partial^{n-2} + \cdots + Q_n(z).
\end{equation}
Then, globally, each coefficient $Q_i$ is in an affine space modeled on
\(
H^0\!\left(C,K_C^i(iD)\right).
\)
\end{itemize}
\end{lemma}

\begin{proposition}\label{p: def twist cotanbun}
Under a change of local coordinate \(z \mapsto w = \phi(z)\) with \(\phi(0)=0\), 
the residue parameters \(v_i\) transform according to
\begin{equation}\label{e:transition for affine bundle}
   v_i^{w}
=
v_i^{z}\,\phi'(0)
-
\frac{n^2+2}{2n}\,
\frac{\phi''(0)}{\phi'(0)} .
\end{equation}
where \(v_i^{w}\) (resp.\ \(v_i^{z}\)) denotes the value of the residue parameter 
in the \(w\)– (resp.\ \(z\)–) trivialization.  
This transformation law defines an affine line bundle \(\mathcal{S}\) modeled on \(T^*C\), 
and each pair \((p_i, v_i)\) determines a point in \(\mathrm{Tot}(\mathcal{S})\).
\end{proposition}

\begin{lemma}\label{l: standard form near div dr}
Let $(\widetilde D,E,\nabla)\in\mathbb{L}_{dR}(d)$. By definition, there exists 
$i: L \hookrightarrow E$ such that 
\[
\widetilde D=\widetilde{D}_i(\nabla)=\{(p_i,v_i)\}_{i=1}^d,
\qquad d=\deg D.
\]
For each $p_i$, choose a local chart $(U_i,z)$ with $z(p_i)=0$.
Then, over $U_i$, with respect to a local frame adapted to the filtration in \eqref{eqn-nabla-filtration}, the dual connection $\nabla^{\vee}|_{U_i}$ is gauge equivalent to
\begin{equation}\label{e: gauge form for nabla Ui}
    \nabla^{\vee} \mid_{U_i} = \mathrm{d} +
    \begin{pmatrix}
        &v^{(z)}_i  &\ast &\ast &\dots &\ast &\ast &\ast  \\
        &-z &-v^{(z)}_i &\ast &\dots &\ast &\ast &\ast \\
        &0 &-1  &0 &\dots &\ast &\ast &\ast \\
        &\vdots &\vdots &\vdots &\dots &\vdots  &\vdots &\vdots \\
        &0  &0 &0 &\dots &-1 &0 &\ast \\
        &0  &0 &0 &\dots &0 &-1 & 0 
    \end{pmatrix} \mathrm{d}z,
\end{equation}
where $v^{(z)}_i$ is the residue parameter under the trivialization over $U_i$ in \Cref{p: def twist cotanbun}, and all unspecified entries are holomorphic in $z$.
\end{lemma}

\begin{theorem}\label{thm: Lag corres dR; SLn}
The variety $\mathbb{L}_{dR}(d)$ defines a \emph{Lagrangian correspondence} between
\(
\mathcal{M}_{dR}(n, \mathcal{O}_C)
\)
and
\(
\operatorname{Hilb}^d(\mathcal{S}).
\)
In other words, $\mathbb{L}_{dR}(d)$ is a Lagrangian subvariety of
\(
\operatorname{Hilb}^d(\mathcal{S})\times \mathcal{M}_{dR}(n, \mathcal{O}_C),
\)
equipped with the holomorphic symplectic form
\(
\omega_{\mathcal{S}}^{[d]}\oplus \Omega_{\mathrm{AB}}.
\)
The same statement holds for $\mathrm{GL}_n$.
\end{theorem}

\begin{proof}[Proof of \Cref{thm: Lag corres dR; SLn}]
By \Cref{p: generic finite dR}, the space $\mathbb{L}_{dR}(d)$ is a half-dimensional subvariety of
\(
\operatorname{Hilb}^d(\mathcal A)\times \mathcal{M}_{dR}(n, \mathcal{O}_C).
\)
It therefore suffices to show that $\mathbb{L}_{dR}(d)$ is isotropic.

Let
\[
(\widetilde D(t_1,t_2),E(t_1,t_2),\nabla(t_1,t_2))
\in \mathbb{L}_{dR}(d)
\]
be a two-parameter deformation family, with $(t_1,t_2)$ sufficiently small.
By \Cref{p: generic finite dR}, this family is in one-to-one correspondence with a two-parameter deformation of triples
\begin{equation}\label{e: two para class dR}
(i(t_1,t_2)\colon L(t_1,t_2)\hookrightarrow E(t_1,t_2),\nabla(t_1,t_2)).
\end{equation}
We first study the deformation of the bundles and connections $(E(t_1,t_2),\nabla(t_1,t_2))$.
Since there is a finite cover from $(E(t_1,t_2),\nabla(t_1,t_2))$ to the associated
$\mathrm{PSL}_n$–flat connection
\(
(\mathbb P(E(t_1,t_2)),\nabla(t_1,t_2)),
\)
it is sufficient to work on the $\mathrm{PSL}_n$ level.

By \Cref{p:trans of diff opers}, once the theta characteristic is fixed, each triple determines a meromorphic operator
\(
\mathbf D_m(t_1,t_2)\in \operatorname{Diff}^n_{C\setminus D} (K_{C\setminus D}^{-\frac{n-1}{2}}, K_{C\setminus D}^{\frac{n+1}{2}}).
\)
Combining this with \Cref{p:oper bundle out D} and the fact that $L(t_1,t_2)^{-1} \cong K_C^{-\frac{n-1}{2}}\otimes \mathcal{O}_C(\frac{D}{n})$, the projective bundle
$\mathbb P(E(t_1,t_2))$ can be obtained as a modification of $\mathbb P(J^{n-1}(K_C^{-\frac{n-1}{2}})^\vee)$ at the divisor
$D(t_1,t_2)=\pi_*(\widetilde D(t_1,t_2))$.

Away from the divisor $D(t_1,t_2)$, we gauge
$\nabla(t_1,t_2)$ (and hence its dual)
into the standard oper form.
More precisely, if $(V,z) \subset C\setminus D(t_1,t_2)$ is a local chart, then
\[
\mathbf D_m(t_1,t_2)\big|_V
=\partial^n+Q_2(z;t_1,t_2)\partial^{n-2}+\cdots+Q_n(z;t_1,t_2),
\]
and the dual connection takes the form
\begin{equation}\label{e: mero oper form dr lag}
\nabla(t_1,t_2)^{\vee}\big|_V
=\mathrm{d}+
\begin{pmatrix}
 0 & Q_2 & \cdots & Q_{n-1} & Q_{n} \\
 -1 & 0 & \cdots & 0 & 0\\
 0 & -1 & \ddots & \vdots & \vdots\\
 \vdots & \vdots & \ddots & 0 & 0\\
 0 & 0 & \cdots & -1 & 0
\end{pmatrix}\mathrm{d}z.
\end{equation}

Now write
\(
D(t_1,t_2)=\sum_i p_i(t_1,t_2)
\).
For each $p_i(0,0)\in C$, fix a neighborhood $(U_i,z)$ with $z(p_i(t_1,t_2))=u_i(t_1,t_2)$, $u_i(0,0)=0$, $\mathbb P(E(t_1,t_2)^{\vee})$ is obtained from $\mathbb P(J^{n-1}(K_C^{-\frac{n-1}{2}}))$ by a transition
function
\[
G_i(z;t_1,t_2)\in
\Gamma(U_i\setminus p_i(t_1,t_2),\mathrm{PSL}_n\otimes\mathcal O),\qquad 1 \leq i \leq d,
\]
satisfying
\begin{equation}\label{e:AS condition ess}
G_i(z;t_1,t_2)\cdot
\nabla(t_1,t_2)^{\vee}\big|_{U_i}
=
\nabla(t_1,t_2)^{\vee}\big|_{U_i\setminus p_i(t_1,t_2)},
\end{equation}
where the RHS is the oper form \eqref{e: mero oper form dr lag}, while LHS is in the form given by \Cref{l: standard form near div dr},
\begin{equation}\label{e:form of nabla check Ui}
\nabla^{\vee}(t_1,t_2)\big|_{U_i}
= \mathrm{d} +
\scalebox{0.85}{$
\begin{pmatrix}
        &v^{(z)}_i(t_1,t_2)  &\ast &\ast &\dots &\ast &\ast &\ast  \\
        &-(z-u_i(t_1,t_2)) &-v^{(z)}_i(t_1,t_2) &\ast &\dots &\ast &\ast &\ast \\
        &0 &-1  &0 &\dots &\ast &\ast &\ast \\
        &\vdots &\vdots &\vdots &\dots &\vdots  &\vdots &\vdots \\
        &0  &0 &0 &\dots &-1 &0 &\ast \\
        &0  &0 &0 &\dots &0 &-1 & 0 
    \end{pmatrix}
$}\,\mathrm{d}z .
\end{equation}
For simplicity, we write $\nabla^{\vee}(t_1,t_2)\big|_{U_i} = \mathrm{d} + A_i(z;t_1,t_2)$ for this gauge form.

Moreover, note that $G_i$ admits a factorization
\[
G_i(z;t_1,t_2)
=(z-u_i(t_1,t_2))^{\omega_1^{\vee}}\cdot N_i(z;t_1,t_2),
\]
with $N_i$ valued in the unipotent lie subgroup $N$.
Here
\[
(z-u_i(t_1,t_2))^{\omega_1^{\vee}}
=
\mathrm{diag}\bigl(
(z-u_i)^{(n-1)/n},
(z-u_i)^{-1/n},\dots,(z-u_i)^{-1/n}
\bigr).
\]
From \eqref{e:AS condition ess},
one finds that $N_i(z;t_1,t_2)$ must be of the form
\begin{equation}\label{e:form of Ni}
N_i(z;t_1,t_2)=
\begin{pmatrix}
 1 & -\frac{n-1}{n(z-u_i(t_1,t_2))^2}+\frac{v^{(z)}_i(t_1,t_2)}{z-u_i(t_1,t_2)} & * & \cdots & * \\
 0 & 1 & -\frac{n-2}{n(z-u_i(t_1,t_2))} & \cdots & * \\
 0 & 0 & 1 & \cdots & * \\
 \vdots & \vdots & \vdots & \ddots & \vdots \\
 0 & 0 & 0 & \cdots & -\frac{1}{n(z-u_i(t_1,t_2))} \\
 0 & 0 & 0 & \cdots & 1
\end{pmatrix}.
\end{equation}

Therefore, evaluating the Atiyah--Bott form $\Omega_{AB}$ on the two
tangent vectors yields,
\begin{align*}
&\Omega_{AB} \left( (E,\nabla)_{t_1}, ( E,\nabla)_{t_2}\right)
= \Omega_{AB} \left( (\mathbb{P}(E),\nabla)_{t_1}, ( \mathbb{P}(E),\nabla)_{t_2}\right) \\
=& 
\sum_{i=1}^d \operatorname{Res}_{z=0} \operatorname{Tr}\left(
    \partial_{t_1} A_i \cdot \partial_{t_2} G_i
    \cdot G_i^{-1}
    - \partial_{t_2} A_i \cdot \partial_{t_1} G_i
    \cdot G_i^{-1}
\right) \Big|_{t_1=t_2=0} \mathrm{d}z \\
=& \sum_{i=1}^d \operatorname{Res}_{\lambda = 0} 
\left[
    -\frac{2(n-1)}{n z^2} \left( \partial_{t_1}u_i + \partial_{t_2}u_i \right)
    + 
    \frac{1}{z} \left( \partial_{t_2} v^{(z)}_i \partial_{t_1} u_i - \partial_{t_1} v^{(z)}_i \partial_{t_2} u_i \right)
    + 
    \mathrm{O}(1)
\right] d\lambda \\
=& \sum_{i=1}^d -\left(
    \partial_{t_1} v^{(z)}_i \cdot \partial_{t_2} u_i
    - \partial_{t_2} v^{(z)}_i \cdot \partial_{t_1} u_i
\right) \Big|_{t_1=t_2=0}\\
=& -\omega^{[d]} \left( \widetilde{D}_{t_1}, \widetilde{D}_{t_2} \right).
\end{align*}
Here $(E,\nabla)_{t_j}$ denotes the tangent vector
\(
(E,\nabla)_{t_j}
:=\partial_{t_j}(E(t_1,t_2),\nabla(t_1,t_2))\Big|_{t_1=t_2=0}.
\)
In the second equality we use \eqref{e:form of nabla check Ui} and
\eqref{e:form of Ni}, from which it follows that only the $2\times2$ minors
contribute to the trace.
In the last line, we write
\(
\widetilde D_{t_j}
:=\partial_{t_j}\widetilde D(t_1,t_2)\Big|_{t_1=t_2=0}.
\)

Therefore, this shows that $\mathbb L_{dR}(d)$ is isotropic, and hence gives a Lagrangian correspondence.
\end{proof}

\end{document}
